\chapter{КОНФИГУРАЦИЯ ЛОКАЛЬНОГО ЭКСПЕРИМЕНТА}
\label{app:configuration}

\begin{table}[H]
    \centering
    \caption{Параметры серии Qwen3-8B-Q4\_K\_M}
    \label{tab:experiment-config}
    \small
    \begin{tabularx}{\textwidth}{L{5cm}X}
        \toprule
        Параметр & Значение \\
        \midrule
        Checkpoint & Qwen/Qwen3-8B-GGUF \\
        Revision & 7c41481f57cb95916b40956ab2f0b139b296d974 \\
        Файл & Qwen3-8B-Q4\_K\_M.gguf, 5 027 783 488 байт \\
        GPU & NVIDIA GeForce RTX 5090 Laptop, 24 463 MiB, WSL2 \\
        Inference & vLLM 0.30.0 и официальный GGUF plugin \\
        Tokenizer/config & Qwen/Qwen3-8B \\
        Контекст & 16 384 токена \\
        Max output & 8 192 токена \\
        Concurrency & 1 \\
        Max repairs & 2 \\
        Pipeline SHA-256 & {\scriptsize\ttfamily 62695337cf3f4912\allowbreak 7909d0cb567fa98e\allowbreak 95d37faf98c9df85\allowbreak 4ec169f851d37f91} \\
        \bottomrule
    \end{tabularx}
\end{table}

Температура задаётся по этапам: 0,1 для brief и components, 0,2 для variables и equations, 0,1 для repair. Дополнительные sampling-параметры берутся из конфигурации модели и vLLM и должны сохраняться вместе с отчётом. Единый seed для побайтовой детерминированности не устанавливался, поэтому повторы не считаются идентичными запусками.

Пик мониторинга составил 21 352 MiB VRAM, 80 °C и 176,91 Вт. После завершения vLLM занято 225 MiB. Значения памяти зависят от выбранной доли `gpu-memory-utilization` и размера KV-cache и не являются минимальными требованиями checkpoint.

\section{Controlled llama.cpp-серии}

Controlled серии 8B и 30B-A3B выполнены последовательно на одном GPU, без одновременного размещения моделей. Использованы одинаковые runtime, source, corpus и inference settings; результаты не смешиваются с прежними acceptance-таблицами.

\begin{table}[H]
    \centering
    \caption{Общие параметры controlled локальных Qwen3}
    \label{tab:controlled-llama-config}
    \small
    \begin{tabularx}{\textwidth}{L{5cm}X}
        \toprule
        Параметр & Значение \\
        \midrule
        Runtime & llama.cpp server-cuda, build b12169-3ed7284f2 \\
        Docker digest & {\scriptsize\ttfamily sha256:5bb779a449a18d8f\allowbreak 4333e9814e88774634\allowbreak 2fe555c528f5eabc92b2\allowbreak 093e32673a} \\
        Pipeline SHA-256 & {\scriptsize\ttfamily cfc09986bc008ec746\allowbreak 28ff54efeff4a1bcc3\allowbreak 19611601c52962f8e0fcd\allowbreak 77c844e} \\
        Контекст / max output & 8192 / 4096 токенов \\
        Sampling & temperature 0,3; top\_k 40; top\_p 0,95; min\_p 0,05 \\
        Inference seeds & $2026+r$, где $r=0,1,2$ --- индекс повтора \\
        Reasoning format & none \\
        Structured output & Полные JSON Schema; ComponentPlan --- JSON-only \\
        Pattern shortcut / RAG & Включён / выключен \\
        Unstructured fallback & Выключен \\
        Request retries / timeout & 2 / 600 с \\
        Concurrency / max repairs & 1 / 2 \\
        \bottomrule
    \end{tabularx}
\end{table}

\begin{table}[H]
    \centering
    \caption{Идентификация controlled checkpoints}
    \label{tab:controlled-checkpoints}
    \small
    \begin{tabularx}{\textwidth}{L{4cm}XX}
        \toprule
        Параметр & 8B & 30B-A3B \\
        \midrule
        Repository & Qwen/Qwen3-8B-GGUF & Qwen/Qwen3-30B-A3B-GGUF \\
        Revision & {\scriptsize\ttfamily 7c41481f57cb9591\allowbreak 6b40956ab2f0b139b296d974} & {\scriptsize\ttfamily e4d4bafdfb96a411\allowbreak a163846265362aceb0b9c63a} \\
        Файл & Qwen3-8B-Q4\_K\_M.gguf & Qwen3-30B-A3B-Q4\_K\_M.gguf \\
        Размер файла, байт & 5 027 783 488 & 18 556 685 824 \\
        Параметров всего & 8 190 735 360 & 30 532 122 624 \\
        Alias & small & large \\
        \bottomrule
    \end{tabularx}
\end{table}

SHA-256 файла 8B:
\begin{center}
\scriptsize\ttfamily d98cdcbd03e17ce476\allowbreak 81435b5150e34c1417f5\allowbreak 0b5c0019dd560e4882c5\allowbreak 745785
\end{center}
SHA-256 файла 30B-A3B:
\begin{center}
\scriptsize\ttfamily 0d003f6662faee786e\allowbreak d5da3e31b29c978de5ae5d\allowbreak 275c8794c606a7f3c01aa8f5
\end{center}
Проверка контрольных сумм выполнена до inference; revision и quantization зафиксированы в манифестах.

Для каждой модели выполнено 216 случаев: 24 prompts, три повтора и три ветви генерации. Пик 8B-мониторинга --- 6536 MiB VRAM, 78 °C и 172,97 Вт; после её серии наблюдалось 513 MiB used и 23 625 MiB free. Пик 30B-A3B --- 19 216 MiB VRAM, 75 °C и 170,87 Вт; после автоматической остановки inference и monitor --- 487 MiB used и 23 651 MiB free. Активных CUDA-процессов после каждой остановки не наблюдалось. Evaluation-контейнеры 30B удалены, дальнейшие GPU-запуски в этой серии не выполняются.
